\section*{Bab \ref{ch:b1:poly} --- Polinomial}

\begin{solution}{exo:b1:poly:1}
$X^5 - 1 = (X^2 + X + 1)(X^3 - X^2 + 1) + (-X - 2)$. Langkahnya: kurangkan
$X^3 B$, lalu $-X^2 B$, lalu $B$; sisanya $-X - 2$ berderajat
$1 < 2$. \emph{Periksa di $X = 1$:} $\;0 = 3 \times 1 + (-3)$.

$2X^4 + X^3 - X + 3 = (X^2 - 2)(2X^2 + X + 4) + (X + 11)$. \emph{Periksa
di $X = 0$:} $\;3 = (-2)(4) + 11$.
\end{solution}

\begin{solution}{exo:b1:poly:2}
$X^2 + X + 1 = (X - j)(X - j^2)$ dengan $j = \eu^{2\iu\pi/3}$, $j^3 =
1$. Ia \omterm{def:b1:arith:divides}{membagi} $Q_n = X^{2n} + X^n + 1$ jika dan hanya jika $j$ dan $j^2$ merupakan akar
$Q_n$; dan karena $Q_n$ berkoefisien real, $Q_n(j^2) =
\conj{Q_n(j)}$, sehingga syaratnya cukup $Q_n(j) = 0$. Sekarang $Q_n(j) =
j^{2n} + j^n + 1$ bergantung pada $n$ modulo $3$:
\begin{itemize}
  \item $n \equiv 0$: $Q_n(j) = 1 + 1 + 1 = 3 \neq 0$;
  \item $n \equiv 1$: $Q_n(j) = j^2 + j + 1 = 0$;
  \item $n \equiv 2$: $Q_n(j) = j^4 + j^2 + 1 = j + j^2 + 1 = 0$.
\end{itemize}
Jadi $X^2 + X + 1 \mid X^{2n} + X^n + 1$ tepat ketika $3 \nmid n$.
\end{solution}

\begin{solution}{exo:b1:poly:3}
Menurut \cref{prop:b1:poly:multiplicity}, $(X-1)^2 \mid P$ jika dan hanya jika $P(1) =
P'(1) = 0$:
\[
P(1) = 2 + a + b = 0, \qquad P'(1) = 4 + 3a + 2b = 0 .
\]
Dengan menyelesaikannya: $b = -a - 2$ dan $4 + 3a - 2a - 4 = a = 0$, jadi $a = 0$, $b =
-2$: sehingga $P = X^4 - 2X^2 + 1 = (X^2 - 1)^2 = (X-1)^2 (X+1)^2$, dan itulah
pemfaktoran realnya.
\end{solution}

\begin{solution}{exo:b1:poly:4}
$X^3 - 1 = (X - 1)(X - j)(X - j^2)$ atas $\C$ (dengan $j = \eu^{2\iu\pi/3}$),
dan $(X - 1)(X^2 + X + 1)$ atas $\R$.

$X^4 + X^2 + 1 = (X^2 + X + 1)(X^2 - X + 1)$ atas $\R$ (jabarkan saja,
atau perhatikan $X^4 + X^2 + 1 = (X^2+1)^2 - X^2$); sedangkan atas $\C$, masing-masing kuadratnya
terbelah: dengan akar $j, j^2$ dan $-j, -j^2$, yakni $\eu^{\pm 2\iu\pi/3},
\eu^{\pm\iu\pi/3}$.

$X^6 - 1 = \prod_{k=0}^{5} (X - \eu^{\iu k\pi/3})$ atas $\C$, dan atas
$\R$:
\[
X^6 - 1 = (X-1)(X+1)(X^2 + X + 1)(X^2 - X + 1),
\]
dengan mengelompokkan pasangan \omterm{def:b1:complex:field}{konjugat} $\eu^{\pm 2\iu\pi/3}$ dan $\eu^{\pm
\iu\pi/3}$.
\end{solution}

\begin{solution}{exo:b1:poly:5}
\begin{enumerate}
  \item Misalkan $p/q$ (dalam bentuk paling sederhana) akar \omterm{def:b1:poly:def}{polinomial
        monik} berkoefisien bulat $X^3 + \dots + c_0$: menghilangkan penyebutnya pada
        $P(p/q) = 0$ memberikan $p^3 = -q\,(\text{bilangan bulat})$, jadi $q \mid
        p^3$; lalu kesalingprimaannya memaksa $q = \pm 1$: sehingga akarnya bilangan bulat
        $p$, dan $p \mid c_0$ (isolasikan $c_0$). Di sini calonnya
        \omterm{def:b1:arith:divides}{membagi} $6$: dengan menguji, $P(1) = 0$, $P(2) = 0$, $P(3) = 0$. Jadi
        $P = (X-1)(X-2)(X-3)$.
  \item Vieta: $s_1 = 6$, $s_2 = 11$, $s_3 = 6$. Jumlah kuadratnya:
        $s_1^2 - 2s_2 = 36 - 22 = 14 = 1 + 4 + 9$, sesuai harapan. Jumlah kebalikannya:
        $\frac{s_2}{s_3} = \frac{11}{6} = 1 + \frac12 + \frac13$, sesuai harapan.
\end{enumerate}
\end{solution}

\begin{solution}{exo:b1:poly:6}
Langkah pembagian pertama pada \omterm{met:b1:arith:euclid}{algoritma Euclid}:
\[
(X + 1)(X^3 - X^2 + X - 1)
= X^4 - X^3 + X^2 - X + X^3 - X^2 + X - 1 = X^4 - 1 ,
\]
jadi pembagian $X^4 - 1$ oleh $X^3 - X^2 + X - 1$ bersifat eksak
(dengan hasil bagi $X + 1$ dan sisa $0$), sehingga algoritmanya langsung berhenti:
\[
\gcd(X^4 - 1,\; X^3 - X^2 + X - 1) = X^3 - X^2 + X - 1
\]
(yang memang sudah \omterm{def:b1:poly:def}{monik}). Hubungan Bézoutnya pun sepele: $\gcd = 0
\cdot (X^4 - 1) + 1 \cdot (X^3 - X^2 + X - 1)$. Pemeriksaan kesejalanannya lewat
pemfaktoran: $X^3 - X^2 + X - 1 = (X - 1)(X^2 + 1)$, yang memang
merupakan hasil kali faktor tak tereduksikan persekutuan $X^4 - 1 =
(X-1)(X+1)(X^2+1)$.
\end{solution}

\begin{solution}{exo:b1:poly:7}
Karena $P \geq 0$ pada $\R$, maka akar realnya berkegandaan genap (karena pada
akar yang berkegandaan ganjil, $P$ berganti tanda). Dengan memakai
\cref{cor:b1:poly:factorization} beserta pemasangannya, tulis
\[
P = c \prod_i (X - a_i)^{2k_i} \prod_j \bigl((X - z_j)(X - \conj
z_j)\bigr)^{n_j},
\]
dengan $c > 0$ (dari perilakunya di $+\infty$). Misalkan
\[
S = \sqrt c\, \prod_i (X - a_i)^{k_i} \prod_j (X - z_j)^{n_j}
\in \C[X],
\]
sehingga $P = S\,\conj S$ dengan $\conj S$ berkoefisien yang
terkonjugatkan. Pilah $S = A + \iu B$ dengan $A, B \in \R[X]$: maka
\[
P = (A + \iu B)(A - \iu B) = A^2 + B^2 .
\]
\end{solution}

\begin{solution}{exo:b1:poly:8}
Lagrange (\cref{thm:b1:poly:lagrange}) dengan simpul $0, 1, 2$:
\[
P = 1\cdot\frac{(X-1)(X-2)}{(0-1)(0-2)} + 3\cdot\frac{X(X-2)}{1\cdot(1-2)}
+ 2\cdot\frac{X(X-1)}{2\cdot 1}
= \frac{(X-1)(X-2)}{2} - 3X(X-2) + X(X-1).
\]
Dengan menjabarkannya: $\frac{X^2 - 3X + 2}{2} - 3X^2 + 6X + X^2 - X =
-\frac{3}{2}X^2 + \frac{7}{2}X + 1$.

Sistemnya: $P = aX^2 + bX + c$ dengan $c = 1$; $a + b + 1 = 3$; $4a + 2b +
1 = 2$. Dengan mengurangkan dua kali persamaan kedua dari yang ketiga: $2a - 1 = -4$,
jadi $a = -\frac32$, $b = \frac72$. \omterm{def:b1:poly:def}{Polinomialnya} sama:
$P = -\frac32 X^2 + \frac72 X + 1$. (Periksa $P(2) = -6 + 7 + 1 = 2$.)
\end{solution}

\begin{solution}{exo:b1:poly:9}
$P = X^{2n+1} - 1$: di sini $P' = (2n+1)X^{2n} \geq 0$, jadi fungsi
\omterm{def:b1:poly:def}{polinomialnya} naik (secara tegas kecuali di $0$), dengan limit
$\mp\infty$: sehingga ia lenyap tepat sekali pada $\R$ (yaitu di $x = 1$).

Misalkan $E_n = \sum_{k=0}^{n} \frac{X^k}{k!}$. Maka $E_n' = E_{n-1} = E_n
- \frac{X^n}{n!}$. Akar ganda $a$ akan memenuhi $E_n(a) =
E_n'(a) = 0$ (\cref{prop:b1:poly:multiplicity}), sehingga
$\frac{a^n}{n!} = E_n(a) - E_n'(a) = 0$, jadi $a = 0$; padahal $E_n(0) = 1
\neq 0$. Jadi tak ada akar ganda.
\end{solution}

\begin{solution}{exo:b1:poly:10}
\begin{enumerate}
  \item Induksi (kedua kasus basisnya berlaku). Dengan memakai $\cos(n+1)\theta +
        \cos(n-1)\theta = 2\cos\theta\cos n\theta$:
        \[
        T_{n+1}(\cos\theta) = 2\cos\theta \cos n\theta -
        \cos(n-1)\theta = \cos(n+1)\theta .
        \]
  \item $T_n(\cos\theta) = 0$ jika dan hanya jika $n\theta \equiv \frac\pi2 \pmod
        \pi$: jadi bilangan
        \[
        x_k = \cos\Bigl(\frac{(2k+1)\pi}{2n}\Bigr),
        \qquad k = 0, 1, \dots, n-1,
        \]
        adalah $n$ titik berbeda pada $\intoo{-1}{1}$ (karena sudutnya terletak di
        $\intoo{0}{\pi}$ yang di sana $\cos$ \omterm{def:b1:logic:inj}{injektif}), yang semuanya akar
        $T_n$; dan karena $\deg T_n = n$ (dari rekursinya, dengan
        koefisien utama $2^{n-1}$ untuk $n \geq 1$, lewat induksi),
        maka itulah \emph{semua} akarnya, masing-masing sederhana.
  \item Untuk $x = \cos\theta \in \intcc{-1}{1}$: $\abs{T_n(x)} =
        \abs{\cos n\theta} \leq 1$, dengan kesamaan jika dan hanya jika $n\theta \equiv
        0 \pmod\pi$, yakni pada $n+1$ titik $y_k =
        \cos\frac{k\pi}{n}$, $k = 0, \dots, n$, yang di sana $T_n(y_k) =
        (-1)^k$. (Ekuiosilasi inilah yang menjadikan $2^{1-n}T_n$ sebagai \omterm{def:b1:poly:def}{polinomial
        monik} berderajat $n$ dengan norma supremum terkecil pada
        $\intcc{-1}{1}$ --- yang dibuktikan pada soal akhir pekan
        bab ini.)
\end{enumerate}
\end{solution}

\begin{solution}{exo:b1:poly:11}
Tulis $P = c\prod_i (X - a_i)^{m_i}$. Aturan hasil kali (yang diperluas ke
beberapa faktor) memberikan
\[
P' = c\sum_{i} m_i (X - a_i)^{m_i - 1} \prod_{k \neq i} (X -
a_k)^{m_k},
\]
lalu dengan membaginya dengan $P$: $\frac{P'}{P} = \sum_i \frac{m_i}{X - a_i}$
(sebagai fungsi rasional, yakni di luar akarnya).

Misalkan $w$ sebuah akar $P'$. Jika $w$ salah satu $a_i$, maka ia terletak di
selubung konveksnya secara sepele. Jika tidak, dengan menghitung nilainya di $w$:
\[
0 = \sum_i \frac{m_i}{w - a_i}
= \sum_i m_i\, \frac{\conj w - \conj a_i}{\abs{w - a_i}^2} .
\]
Dengan mengonjugatkannya: $\sum_i \lambda_i (w - a_i) = 0$ dengan $\lambda_i =
\frac{m_i}{\abs{w - a_i}^2} > 0$. Jadi
\[
w = \frac{\sum_i \lambda_i a_i}{\sum_i \lambda_i} :
\]
yaitu kombinasi cembung (dengan bobot positif yang berjumlah $1$ setelah
dinormalkan) atas akar $a_i$. Jadi setiap akar $P'$ terletak di
selubung cembung akar $P$.
\end{solution}

\begin{solution}{exo:b1:poly:12}
Jumlahkan nilai $(1 + X)^n$ pada ketiga \omterm{def:b1:complex:unity}{akar satuan} pangkat
tiga:
\[
2^n + (1 + j)^n + (1 + j^2)^n
= \sum_{k=0}^n \binom nk\,\bigl(1 + j^k + j^{2k}\bigr)
= 3\sum_{k\,:\,3\mid k}\binom nk ,
\]
karena $1 + j^k + j^{2k}$ adalah jumlah geometri yang bernilai $3$ ketika $3
\mid k$ dan bernilai $\frac{j^{3k} - 1}{j^k - 1} = 0$ bila tidak. Sekarang $1 +
j = \frac12 + \iu\frac{\sqrt3}2 = \eu^{\iu\pi/3}$ dan $1 + j^2 =
\conj{1 + j} = \eu^{-\iu\pi/3}$, sehingga $(1+j)^n + (1+j^2)^n =
2\cos\frac{n\pi}3$ dan
\[
\sum_{k\geq0}\binom n{3k} = \frac{2^n + 2\cos\frac{n\pi}3}{3} .
\]
Pemeriksaannya: untuk $n = 3$: $\frac{8 + 2\cos\pi}3 = 2 = \binom30 + \binom33$;
untuk $n = 6$: $\frac{64 + 2}3 = 22 = 1 + 20 + 1$.
\end{solution}

\begin{solution}{pb:b1:poly:1}
\textbf{1.} $T_2 = 2X^2 - 1$; $T_3 = 2X(2X^2 - 1) - X = 4X^3 -
3X$; $T_4 = 2X\,T_3 - T_2 = 8X^4 - 8X^2 + 1$; $T_5 = 2X\,T_4 - T_3
= 16X^5 - 20X^3 + 5X$. Kesamaan $T_3(\cos\theta) = \cos3\theta$
persis sama dengan $\cos3\theta = 4\cos^3\theta - 3\cos\theta$ dari
\cref{ex:b1:complex:cos3}.

\textbf{2.} Benar untuk $n = 1, 2$. Jika $T_{n-1}$, $T_n$ berderajat
$n-1$, $n$ dengan koefisien utama $2^{n-2}$, $2^{n-1}$, maka
$2X\,T_n$ berderajat $n+1$ dengan koefisien utama $2^n$, sedangkan
$T_{n-1}$ berderajat lebih rendah: jadi $T_{n+1}$ berderajat $n + 1$ dengan koefisien
utama $2^n$. Paritasnya: jika $T_{n-1}$ berparitas $n - 1$
dan $T_n$ berparitas $n$, maka $2X\,T_n$ dan $T_{n-1}$ sama-sama berparitas
$n + 1$, sehingga demikian pula $T_{n+1}$.

\textbf{3.} Jika $P(\cos\theta) = \cos n\theta$ untuk setiap $\theta$,
maka $P$ dan $T_n$ bersesuaian di setiap titik $\intcc{-1}1$ --- yaitu
\omterm{def:b1:logic:sets}{himpunan} tak hingga --- sehingga $P - T_n$ mempunyai tak hingga banyak akar dan menjadi
\omterm{def:b1:poly:def}{polinomial} nol (\cref{cor:b1:poly:nroots}).

\textbf{4.} Untuk $x = \cos\theta$: $T_m(T_n(\cos\theta)) =
T_m(\cos n\theta) = \cos(mn\theta) = T_{mn}(\cos\theta)$, dan
$2T_mT_n(\cos\theta) = 2\cos m\theta\cos n\theta = \cos(m+n)\theta
+ \cos\abs{m - n}\theta$. Kedua kesamaannya berlaku pada $\intcc{-1}1$,
sehingga berlaku sebagai kesamaan \omterm{def:b1:poly:def}{polinomial} menurut argumen pertanyaan 3.

\textbf{5.} Di sini $x_k$ adalah $n$ akar sederhana yang berbeda dan
koefisien utamanya $2^{n-1}$:
\[
T_n = 2^{n-1}\prod_{k=0}^{n-1}
\Bigl(X - \cos\frac{(2k+1)\pi}{2n}\Bigr) .
\]
Keberselang-selingannya: sudut $0 < \frac{\pi}{2n} < \frac\pi n <
\frac{3\pi}{2n} < \frac{2\pi}n < \dots < \pi$ berganti-ganti antara
sudut $y$ yaitu $\frac{k\pi}n$ dan sudut $x$ yaitu
$\frac{(2k+1)\pi}{2n}$; dan karena $\cos$ turun tegas pada
$\intcc0\pi$, nilainya berselang-seling dalam urutan terbalik: $y_n <
x_{n-1} < y_{n-1} < \dots < x_0 < y_0$. Di antara dua ekstremum
yang berurutan duduk tepat satu akar, seperti disiratkan gambar $\cos n\theta$.

\textbf{6.} Induksi dengan $2\cosh a\cosh b = \cosh(a + b) +
\cosh(a - b)$ (\cref{prop:b1:functions:hyprules}): $T_{n+1}(\cosh
t) = 2\cosh t\cosh nt - \cosh(n-1)t = \cosh(n+1)t$. Untuk $x \geq
1$, tulis $x = \cosh t$ dengan $t \geq 0$; maka $\eu^t = x +
\sqrt{x^2 - 1}$ dan $\eu^{-t} = x - \sqrt{x^2 - 1}$, sehingga
\[
T_n(x) = \cosh(nt)
= \frac{(x + \sqrt{x^2-1})^n + (x - \sqrt{x^2-1})^n}2 .
\]
Untuk $x > 1$ suku pertamanya melampaui $\frac12(1)^n$ secara tegas dan
tumbuh secara geometri: jadi $T_n(x) > 1$.

\textbf{7.} De Moivre: $\cos n\theta = \Re\bigl((\cos\theta +
\iu\sin\theta)^n\bigr) = \sum_{2j \leq n}\binom n{2j}
\cos^{n-2j}\theta\,(\iu\sin\theta)^{2j}$, dan $(\iu\sin\theta)^{2j}
= (-\sin^2\theta)^j = (\cos^2\theta - 1)^j$. Dengan mensubstitusikan $x =
\cos\theta$ lalu memanggil pertanyaan 3:
\[
T_n(x) = \sum_{0\leq 2j\leq n}\binom n{2j}x^{n-2j}(x^2 - 1)^j .
\]
Untuk $n = 3$: $\binom30 x^3 + \binom32 x(x^2 - 1) = x^3 + 3x^3 -
3x = 4x^3 - 3x$, seperti pada pertanyaan 1.

\textbf{8.} $T_n(1) = \cos(n\cdot0) = 1$; $T_n(-1) = \cos(n\pi) =
(-1)^n$; dan $T_n(0) = \cos\frac{n\pi}2$, yang bernilai $0$ untuk $n$ ganjil dan
$(-1)^{n/2}$ untuk $n$ genap.

\textbf{9.} Induksi untuk $U_n(\cos\theta) =
\frac{\sin(n+1)\theta}{\sin\theta}$: benar untuk $U_0 = 1$ dan $U_1 =
2X$ (karena $\sin2\theta = 2\sin\theta\cos\theta$); dan langkahnya adalah
kesamaan jumlah-ke-hasil-kali $\sin(n+2)\theta = 2\cos\theta\,
\sin(n+1)\theta - \sin n\theta$. Sekarang turunkan
$T_n(\cos\theta) = \cos n\theta$ terhadap $\theta$:
$-\sin\theta\,T_n'(\cos\theta) = -n\sin n\theta$, jadi untuk $\theta
\notin \pi\Z$:
\[
T_n'(\cos\theta) = n\,\frac{\sin n\theta}{\sin\theta}
= n\,U_{n-1}(\cos\theta) ,
\]
dan \omterm{def:b1:poly:def}{polinomial} $T_n'$ serta $nU_{n-1}$, yang bersesuaian pada
$\intoo{-1}1$, adalah sama.

\textbf{10.} $\abs{\sin(n+1)\theta} = \abs{\sin n\theta\cos\theta
+ \cos n\theta\sin\theta} \leq \abs{\sin n\theta} +
\abs{\sin\theta}$, lalu induksi memberikan $\abs{\sin n\theta} \leq
n\abs{\sin\theta}$. Jadi $\abs{U_{n-1}} \leq n$ pada $\intoo{-1}1$
dan $\abs{T_n'} = n\abs{U_{n-1}} \leq n^2$ di sana; sedangkan di $\pm1$
batasnya diperluas lewat limit (atau langsung: $U_{n-1}(1) = n$ dari
rekursinya, $U_n(1) = n + 1$ lewat induksi, dan paritasnya memberikan
$U_{n-1}(-1) = (-1)^{n-1}n$). Jadi $T_n'(1) = n^2$ dan $T_n'(-1)
= (-1)^{n-1}n^2$: sehingga batas $n^2$ tercapai pada titik ujungnya.

\textbf{11.} Turunkan $\sin\theta\,T_n'(\cos\theta) = n\sin
n\theta$ (pertanyaan 9) terhadap $\theta$:
\[
\cos\theta\,T_n'(\cos\theta) - \sin^2\theta\,T_n''(\cos\theta)
= n^2\cos n\theta = n^2\,T_n(\cos\theta) .
\]
Dengan $x = \cos\theta$ dan $\sin^2\theta = 1 - x^2$: $x\,T_n' - (1
- x^2)T_n'' = n^2T_n$ pada $\intcc{-1}1$, sehingga di mana-mana:
jadi $(1 - x^2)y'' - xy' + n^2y = 0$ untuk $y = T_n$. Periksa untuk $T_2 =
2x^2 - 1$: $(1 - x^2)(4) - x(4x) + 4(2x^2 - 1) = 4 - 4x^2 - 4x^2
+ 8x^2 - 4 = 0$.

\textbf{12.} Di sini $\widetilde T_n$ \omterm{def:b1:poly:def}{monik} (pertanyaan 2) dan
$\abs{\widetilde T_n} = 2^{1-n}\abs{T_n} \leq 2^{1-n}$ pada
$\intcc{-1}1$, dengan $\widetilde T_n(y_k) = (-1)^k2^{1-n}$ pada
$n + 1$ titik $y_k$ (\cref{exo:b1:poly:10}): jadi normanya tepat
$2^{1-n}$, yang tercapai dengan tanda berselang-seling.

\textbf{13.} Di sini $\widetilde T_n$ dan $P$ sama-sama \omterm{def:b1:poly:def}{monik} berderajat
$n$, sehingga suku utamanya saling hapus: jadi $\deg D \leq n - 1$. Pada $y_k$:
$D(y_k) = (-1)^k2^{1-n} - P(y_k)$, dan $\abs{P(y_k)} \leq \norm
P_\infty < 2^{1-n}$ memaksa tanda $D(y_k)$ sama dengan tanda
$(-1)^k2^{1-n}$, secara tegas.

\textbf{14.} Di sini $D$ berganti tanda antara $y_{k+1}$ dan $y_k$ untuk masing-masing
$k = 0, \dots, n-1$: jadi menurut sifat nilai antara, $D$ mempunyai
sebuah akar pada masing-masing dari $n$ selang terbuka yang saling lepas itu --- yakni $n$
akar berbeda bagi \omterm{def:b1:poly:def}{polinomial} tak nol berderajat $\leq n - 1$,
yang mustahil. Dan $D = 0$ pun mustahil (karena normanya berbeda).
Jadi kontradiksi: tak ada $P$ \omterm{def:b1:poly:def}{monik} berderajat $n$ dengan $\norm P_\infty <
2^{1-n}$, dan itulah teorema Chebyshev.

\textbf{15.} Sekarang $\abs{P(y_k)} \leq 2^{1-n}$ saja, jadi $(-1)^k
D(y_k) = 2^{1-n} - (-1)^kP(y_k) \geq 2^{1-n} - \abs{P(y_k)} \geq
0$. Andaikan $D(y_k) = 0$ pada sebuah $y_k$ yang di dalam ($0 < k < n$): maka
$P(y_k) = (-1)^k2^{1-n}$, sehingga $\abs P$ mencapai supremumnya
$2^{1-n}$ pada titik dalam $y_k$, yang mengakibatkan $P'(y_k) = 0$
(karena ekstremum di dalam); dan $T_n'(y_k) = nU_{n-1}(y_k) = 0$ karena
$\sin(n\cdot\frac{k\pi}n) = 0$ --- jadi $\widetilde T_n'(y_k) = 0$
juga, sehingga $D'(y_k) = 0$: yakni $y_k$ akar $D$ berkegandaan
sekurang-kurangnya $2$.

\textbf{16.} Cacahlah akar $D$ beserta \omterm{def:b1:poly:derivative}{kegandaannya}. Misalkan $z$
banyaknya titik dalam $y_k$ dengan $D(y_k) = 0$ (masing-masing akar
ganda, menurut pertanyaan 15) dan $e \in \{0, 1, 2\}$ banyaknya
titik ujung ($y_0$ atau $y_n$) dengan $D = 0$ (masing-masing sekurang-kurangnya akar
sederhana). Sebuah celah $(y_{k+1}, y_k)$ yang kedua ujungnya sama-sama mempunyai $D
\neq 0$ memikul tanda yang berselang-seling secara tegas, sehingga ia memuat
akar di dalamnya. Setiap titik dalam yang lenyap merusak paling banyak dua
celah yang bersebelahan dengannya, dan setiap titik ujung yang lenyap paling banyak satu celah: jadi sekurang-kurangnya
$n - 2z - e$ celah tetap menyumbang satu akar masing-masing, yang semuanya berbeda
dari akar $y$ tadi. Totalnya: sekurang-kurangnya $(n - 2z - e) + 2z + e = n$
akar beserta \omterm{def:b1:poly:derivative}{kegandaannya}, bagi \omterm{def:b1:poly:def}{polinomial} berderajat $\leq n - 1$:
jadi $D = 0$ dan $P = \widetilde T_n$. Peminimumnya pun tunggal.

\textbf{17.} \omterm{def:b1:logic:map}{Pemetaan} afin $t \mapsto x = \frac{a+b}2 +
\frac{b-a}2\,t$ adalah bijeksi $\intcc{-1}1 \to \intcc ab$. Jika $P$
\omterm{def:b1:poly:def}{monik} berderajat $n$, maka $Q(t) = P(x(t))$ adalah \omterm{def:b1:poly:def}{polinomial} dalam
$t$ dengan koefisien utama $\bigl(\frac{b-a}2\bigr)^n$, dan
$\sup_{\intcc ab}\abs P = \sup_{\intcc{-1}1}\abs Q$. Adapun \omterm{def:b1:poly:def}{polinomial
monik} $Q/\bigl(\frac{b-a}2\bigr)^n$ bernorma supremum $\geq
2^{1-n}$ (pertanyaan 13--14), sehingga
\[
\sup_{\intcc ab}\abs P \geq \Bigl(\frac{b-a}2\Bigr)^n 2^{1-n}
= 2\Bigl(\frac{b-a}4\Bigr)^n ,
\]
dengan kesamaan tepat untuk $P(x) = \bigl(\frac{b-a}2\bigr)^n
\widetilde T_n\bigl(t(x)\bigr)$ (pertanyaan 16).

\textbf{18.} $\widetilde T_3 = \frac{T_3}4 = X^3 - \frac34X$;
lalu $\widetilde T_3{}' = 3X^2 - \frac34$ lenyap di $\pm\frac12$.
Nilainya: $\widetilde T_3(-1) = -\frac14$, $\widetilde
T_3(-\tfrac12) = \frac14$, $\widetilde T_3(\tfrac12) = -\frac14$,
$\widetilde T_3(1) = \frac14$: yaitu empat ekstremum berselang-seling
bernilai mutlak $\frac14$ --- jadi $\norm{\widetilde T_3}_\infty =
\frac14$, dan menurut teorema Chebyshev tak ada kubik \omterm{def:b1:poly:def}{monik} yang bernorma
supremum lebih kecil pada $\intcc{-1}1$.

\textbf{19.} Di sini $\omega$ \omterm{def:b1:poly:def}{monik} berderajat $n + 1$, sehingga
$\norm\omega_\infty \geq 2^{-n}$ menurut teorema Chebyshev (untuk derajat
$n+1$), dengan kesamaan jika dan hanya jika $\omega = \widetilde T_{n+1}
= 2^{-n}T_{n+1}$ (pertanyaan 16), yakni jika dan hanya jika simpulnya
adalah $n + 1$ akar $T_{n+1}$. Dengan simpul Chebyshev, faktor galatnya
$\norm\omega_\infty$ bernilai $2^{-n}$ --- yaitu yang sekecil mungkin.

\textbf{20.} Karena $5 \times 36^\circ = 180^\circ$, maka $T_5(c) =
\cos180^\circ = -1$: jadi $16c^5 - 20c^3 + 5c + 1 = 0$. Dengan menguji $x =
-1$: $-16 + 20 - 5 + 1 = 0$, dan menjabarkannya membenarkan
\[
16x^5 - 20x^3 + 5x + 1 = (x + 1)\bigl(4x^2 - 2x - 1\bigr)^2 .
\]
Karena $c = \cos36^\circ \neq -1$, maka $c$ merupakan akar $4x^2 - 2x -
1$, yang akarnya $\frac{1 \pm \sqrt5}4$; dan karena $c > 0$,
\[
\cos36^\circ = \frac{1 + \sqrt5}4 .
\]
Kesejalanannya: $\cos72^\circ = T_2(c) = 2c^2 - 1 = 2\cdot\frac{3 +
\sqrt5}8 - 1 = \frac{\sqrt5 - 1}4$, yaitu nilai yang ditemukan pada
\cref{exo:b1:complex:8}.

\textbf{21.} $\sqrt{1.1^2 - 1} = \sqrt{0.21} \approx 0.458$, jadi
$x + \sqrt{x^2-1} \approx 1.558$ dan $(1.558)^{10} \approx 84.5$,
sedangkan $(1.1 - 0.458)^{10} \approx 0.01$: sehingga $T_{10}(1.1) \approx
\frac{84.5 + 0.01}2 \approx 42$. Jadi \omterm{def:b1:poly:def}{polinomial} yang terkurung di
$\intcc{-1}1$ pada selang itu sudah tumbuh melampaui $40$ hanya
sepersepuluh di luar tepinya: keterbatasan pada sebuah ruas tak mengatakan apa-apa
sejengkal di luarnya.

\textbf{22.} Pada rumus $T_p$ di pertanyaan 7, suku $j = 0$
adalah $X^p$; sedangkan setiap suku lainnya memikul $\binom p{2j}$ dengan $0 < 2j <
p$ (perhatikan $2j \neq p$ karena $p$ ganjil), yang habis dibagi $p$
menurut langkah pertama bukti \cref{thm:b1:arith:fermat}.
Jadi setiap koefisien $T_p - X^p$ merupakan kelipatan $p$.
Pemeriksaannya: $T_3 - X^3 = 3X^3 - 3X = 3(X^3 - X)$; $T_5 - X^5 = 15X^5
- 20X^3 + 5X = 5(3X^5 - 4X^3 + X)$.

\textbf{23.} Menurut pertanyaan 17 dengan $\intcc ab = \intcc01$ dan $n =
2$: simpangan minimalnya $2\bigl(\frac14\bigr)^2 = \frac18$,
yang tercapai oleh $\bigl(\frac12\bigr)^2\widetilde T_2(2x - 1) =
\frac14\bigl((2x-1)^2 - \frac12\bigr) = x^2 - x + \frac18$. Jadi
kuadrat \omterm{def:b1:poly:def}{monik} yang paling dekat dengan nol pada $\intcc01$ adalah $x^2 - x +
\frac18$, dengan norma supremum $\frac18$.

\textbf{24.} (i) Kekakuannya --- bahwa \omterm{def:b1:poly:def}{polinomial} dengan akar lebih banyak daripada
derajatnya adalah nol --- menggerakkan asas ketunggalannya
(pertanyaan 3), pemindahan kesamaan trigonometri menjadi
kesamaan \omterm{def:b1:poly:def}{polinomial} (pertanyaan 4, 7, 9, 11), dan kedua
argumen pencacahan akar pada bukti keekstremannya (pertanyaan 14,
16). (ii) Trigonometri \cref{ch:b1:complex} (de Moivre dan
jumlah-ke-hasil-kali) beserta \omterm{def:b1:functions:hyperbolic}{fungsi hiperbolik} pada
\cref{ch:b1:functions} memasok setiap kesamaan di balik keluarga itu;
dan substitusi $x = \cos\theta$ menjadi jembatannya. (iii) Adapun
\omterm{def:b1:arith:divides}{keterbagian} $p \mid \binom p{2j}$ dari \cref{ch:b1:arith} mengubah
rumus koefisiennya menjadi \omterm{def:b1:arith:congruence}{kekongruenan} pertanyaan 22.

\textbf{25.} Teorema Chebyshev mengubah pengoptimuman atas sebuah
keluarga berdimensi tak hingga (yaitu semua \omterm{def:b1:poly:def}{polinomial monik}) menjadi
kombinatorika yang hingga: karena pesaing yang lebih baik daripada $\widetilde T_n$ akan
berselisih darinya sebesar \omterm{def:b1:poly:def}{polinomial} berderajat rendah yang terpaksa berganti tanda
$n$ kali --- yakni satu akar lebih banyak daripada yang diizinkan derajatnya. Jadi pola
ekuiosilasinya bukan keanehan melainkan justru
sertifikat keoptimalannya, dan kasus kesamaannya mempertajam
pencacahan akar dengan \omterm{def:b1:poly:derivative}{kegandaannya}. Adapun substitusi $x =
\cos\theta$ pantas mendapat kata terakhir: ia mengangkut dunia \omterm{def:b1:poly:def}{polinomial}
yang kaku dan diskret ke dunia periodik
trigonometri, yang di sana akar dan ekstremum $T_n$ hanyalah
kisi teratur $\cos n\theta$. Kedua fakta analisis yang dipinjam
--- yaitu sifat nilai antara (pertanyaan 14; dibuktikan pada
\cref{ch:b1:continuity}) dan lenyapnya turunan pada
ekstremum di dalam (pertanyaan 15; dibuktikan pada \cref{ch:b1:derivative})
--- justru perkakas yang akan dikembalikan bab berikutnya,
sehingga lingkarannya tertutup.

\end{solution}
